\section{Construcción de redes neuronales con perceptrones}

\subsection{Perceptrones de salida múltiple y capas Dense}

Un solo perceptrón tiene una sola salida. Al colocar varios perceptrones en paralelo, con todas las entradas conectadas simultáneamente a cada perceptrón, se forma una \textbf{capa Dense} (capa totalmente conectada).

\begin{figure}[H]
\centering
\includegraphics[width=0.85\textwidth]{slides-images/slide-036.jpg}
\caption{Perceptrón de salida múltiple: todas las entradas conectadas densamente a todas las salidas, es decir, una capa Dense\protect\footnotemark}
\end{figure}
\footnotetext{Diapositiva 36 de las láminas.}

Para la $i$-ésima neurona de salida:
$$z_i = w_{0,i} + \sum_{j=1}^{m} x_j \cdot w_{j,i}$$

Debido a que todas las entradas están conectadas a todas las salidas, este tipo de capa se llama capa \textbf{Dense} (densamente conectada) o \textbf{capa totalmente conectada} (Fully Connected Layer).

Implementar una capa Dense en código es muy sencillo:

\begin{lstlisting}[caption={Implementación de la capa Dense en TensorFlow y PyTorch}]
# TensorFlow
import tensorflow as tf
layer = tf.keras.layers.Dense(units=2)

# PyTorch
import torch.nn as nn
layer = nn.Linear(in_features=m, out_features=2)
\end{lstlisting}

\begin{figure}[H]
\centering
\includegraphics[width=0.95\textwidth]{slides-images/slide-037.jpg}
\caption{Implementación de la capa Dense desde cero: comparación de código entre TensorFlow (izquierda) y PyTorch (derecha)\protect\footnotemark}
\end{figure}
\footnotetext{Diapositiva 37 de las láminas.}

\subsection{Red neuronal de una sola capa oculta}

Al agregar una capa Dense entre la capa de entrada y la capa de salida como \textbf{capa oculta} (Hidden Layer), se forma una red neuronal de una sola capa oculta:

\begin{figure}[H]
\centering
\includegraphics[width=0.85\textwidth]{slides-images/slide-039.jpg}
\caption{Estructura de una red neuronal de una sola capa oculta\protect\footnotemark}
\end{figure}
\footnotetext{Diapositiva 39 de las láminas.}

Cálculo de la capa oculta:
$$z_i = w_{0,i}^{(1)} + \sum_{j=1}^{m} x_j \cdot w_{j,i}^{(1)}$$

Cálculo de la capa de salida (usando como entrada los valores de activación de la capa oculta):
$$\hat{y}_i = g\left(w_{0,i}^{(2)} + \sum_{j=1}^{d_1} g(z_j) \cdot w_{j,i}^{(2)}\right)$$

Donde los superíndices $(1)$ y $(2)$ indican, respectivamente, los parámetros de la primera y la segunda capa, y $d_1$ es el número de neuronas de la capa oculta.

\begin{importantbox}{El significado de la capa oculta}
La capa oculta se llama «oculta» porque, en los datos de entrenamiento, solo podemos observar directamente la entrada y la salida, mientras que la representación de la capa oculta es aprendida automáticamente por la red y resulta «invisible» para un observador externo. Cada neurona de la capa oculta aprende alguna transformación no lineal de la entrada, y estas transformaciones en conjunto aportan una representación de características enriquecida para la predicción final.
\end{importantbox}

\subsection{Redes neuronales profundas}

Al apilar más capas ocultas, se forma una \textbf{red neuronal profunda} (Deep Neural Network). De ahí proviene precisamente el «profunda» en «aprendizaje profundo».

\begin{figure}[H]
\centering
\includegraphics[width=0.85\textwidth]{slides-images/slide-042.jpg}
\caption{Red neuronal profunda: apilamiento de múltiples capas ocultas\protect\footnotemark}
\end{figure}
\footnotetext{Diapositiva 42 de las láminas.}

La fórmula de cálculo para la $i$-ésima neurona de la capa $k$ es:
$$z_{k,i} = w_{0,i}^{(k)} + \sum_{j=1}^{n_{k-1}} g(z_{k-1,j}) \cdot w_{j,i}^{(k)}$$

Donde $n_{k-1}$ es el número de neuronas de la capa $k-1$.

Construir una red profunda en código resulta igual de directo:

\begin{lstlisting}[caption={Redes profundas en TensorFlow y PyTorch}]
# TensorFlow
model = tf.keras.Sequential([
    tf.keras.layers.Dense(n1),
    tf.keras.layers.Dense(n2),
    # ...
    tf.keras.layers.Dense(2)
])

# PyTorch
model = nn.Sequential(
    nn.Linear(m, n1), nn.ReLU(),
    nn.Linear(n1, n2), nn.ReLU(),
    # ...
    nn.Linear(nK, 2)
)
\end{lstlisting}

\subsection{Resumen de la sección}

Partiendo de un solo perceptrón, al colocar varios perceptrones en paralelo se forma una capa Dense, y al apilar varias capas Dense se forma una red neuronal profunda. La salida de cada capa sirve de entrada a la siguiente, extrayendo capa por capa características cada vez más abstractas. La «profundidad» de una red neuronal profunda proviene de esta estructura multinivel.

